\documentclass[a4paper,10pt]{article}
\usepackage[margin=22mm]{geometry}
\usepackage[T1]{fontenc}
\usepackage[utf8]{inputenc}
\usepackage{microtype,parskip,booktabs,array}
\usepackage[colorlinks=true,allcolors=blue,pdfauthor={}]{hyperref}
\pagestyle{plain}
\begin{document}
\begin{center}
{\Large\bfseries GDTCR: journal fit and submission route}\\[4pt]
Official instructions and 2025 impact factors checked on 3 November 2026
\end{center}

\textbf{Recommendation.}
Prepare the current study as a BMC Bioinformatics research article.
Briefings in Bioinformatics is a conditional alternative if the supervisor
considers the comparative insights sufficiently broad. Bioinformatics is a
weaker fit for the present contribution because its instructions emphasize
substantial methodological advances. Impact factor alone does not determine
fit, and this assessment does not predict acceptance.

\small
\begin{tabular}{>{\raggedright\arraybackslash}p{0.23\linewidth}
                >{\raggedright\arraybackslash}p{0.33\linewidth}
                >{\raggedright\arraybackslash}p{0.36\linewidth}}
\toprule
Journal / 2025 JIF & Relevant format & Assessment for this study\\
\midrule
BMC Bioinformatics / 4.4 &
Research article. Abstract at most 350 words: Background, Results,
Conclusions; 3--10 keywords. Main text uses Background, Methods, Results,
Discussion and Conclusions. &
Best current fit: sequence analysis, computational modeling and
scientific validity. Preserve controls, negative findings and bounded
claims; finalize declarations, access and licensing before submission.\\[7pt]
Briefings in Bioinformatics / 7.3 &
Methods comparison/benchmarking is within scope. Problem-solving protocols
and case studies: 2000--5000 words; the selected article category must fit
the contribution. &
Possible higher-impact alternative. The single model family and
source-concentrated corpus limit breadth. A broader method comparison and
a sufficiently independent cohort would strengthen the case; these are
study-specific recommendations, not universal journal requirements.\\[7pt]
Bioinformatics / 5.5 &
Original Papers: up to 7 formatted pages, approximately 5000 words excluding
figures. New methods must be compared with state-of-the-art methods. &
Current priority is lower. The journal discourages routine applications and
incremental benchmark gains without a substantial methodological advance.
Additional formatting alone would not address that contribution gap.\\
\bottomrule
\end{tabular}
\normalsize

\textbf{Practical route.}
Review the current scientific draft first. If BMC remains the target, finalize
authorship, statements, permissions, code/weight licenses and the data/code
version, then apply the final journal checklist. A new external cohort,
broader modern baselines or matched-total-computation experiments should be
treated as separately scoped extensions if the supervisor requests them;
they are not claimed as completed in this draft.

\textbf{Official sources.}
The figures above are the same-year journal impact factors, not CiteScores or
five-year impact factors.
\begin{enumerate}
\setlength{\itemsep}{1pt}
\item BMC:
\href{https://link.springer.com/journal/12859}{metrics},
\href{https://link.springer.com/journal/12859/aims-and-scope}{scope},
\href{https://link.springer.com/journal/12859/submission-guidelines/research-article}{research-article instructions}.
\item Briefings:
\href{https://academic.oup.com/bib/pages/About}{metrics},
\href{https://academic.oup.com/bib/pages/General_Instructions}{article categories and scope}.
\item Bioinformatics:
\href{https://academic.oup.com/bioinformatics/pages/About}{metrics},
\href{https://academic.oup.com/bioinformatics/pages/General_Instructions}{article categories and contribution criteria}.
\end{enumerate}
Retrieved excerpts, access dates and page hashes are retained in
\texttt{source\_checks/}. No editor has been contacted and no submission
decision has been made.
\end{document}
